\documentclass[12pt]{article}
\usepackage[utf8]{inputenc}
\usepackage[T1]{fontenc}
\usepackage{hyperref}
\title{Moji Language Resources for K-12 Curriculum Development}
\author{Compiled by Claude Sonnet 5.5 (claude-sonnet-5-5) for Andr\'as Kornai}
\date{}
\begin{document}
\maketitle

\section{Introduction}
Moji (ISO 639-3: ymi; also Muji or Pingtou Phula) is a Loloish language of the Sino-Tibetan family, spoken by the Phula in Wenshan and Xichou counties of Yunnan, China. Wikipedia gives about 2,000 speakers (2011). Verbix gives more than 400 and rates the language as vulnerable. Joshua Project's figure of 7,100 is a population count for the people group, not a speaker count. The Moji are officially classified within the Yi nationality, and younger people are shifting toward Mandarin. No educational materials were found.

\section{Primary Educational Resources}
No primers, textbooks or school programmes in Moji were found in the searches.

\section{Linguistic and Reference Materials}
\begin{itemize}
\item Wikipedia, Moji language: \url{https://en.wikipedia.org/wiki/Moji_language}
\item Wikipedia, Loloish languages: \url{https://en.wikipedia.org/wiki/Loloish_languages}
\item Wikipedia, Muzi language (a related language): \url{https://en.wikipedia.org/wiki/Muzi_language}
\item Verbix, Endangered Languages: Moji: \url{https://docs.verbix.com/EndangeredLanguages/Moji}
\item Endangered Languages Project entry: \url{https://endangeredlanguages.com/lang/3179?hl=zh-cn}
\end{itemize}

\section{Religious and Community Resources}
\begin{itemize}
\item Joshua Project, Moji in China: \url{https://joshuaproject.net/people_groups/20605/CH}
\item Joshua Project, language page: \url{https://joshuaproject.net/languages/ymi}
\end{itemize}

\section{Digital and Audio Resources}
No audio, video or software resources specific to Moji were found. The Endangered Languages Project entry may link to further material, but this was not checked.

\section{Curriculum Development Context}
Little is documented and the speaker base is small. No school use or orthography was found. Any curriculum would first need documentation work and community involvement.

\end{document}